\setcounter{chapter}{-1}
\chapter{Guideline}\label{ch:guideline}
\markboth{Guideline}{}

\section{Format and Style}

\subsection{Page Layout and Typography}
\begin{itemize}
    \item Paper format: white DIN A4 paper
    \item Font: Arial, 12pt (in this LaTeX template, the sans-serif option achieves this)
    \item Line spacing: 1.2 (the template uses 1.25)
    \item Text alignment: justified (Blocksatz) with automatic hyphenation enabled
    \item Font emphasis: use only in exceptional cases; headings follow the template style and must not be underlined
    \item Bullet points: only square bullet points (FAPS green), never round ones; use the template's built-in list style
    \item FAPS color scheme: FAPS-Grün (RGB 151, 193, 57) and white
\end{itemize}

\subsection{Page Numbering and Binding}
\begin{itemize}
    \item Page numbering starts with the introduction (Arabic numerals)
    \item Directories (table of contents, list of figures, etc.) receive separate Roman numeral numbering; the cover page is not counted
    \item Print single-sided unless otherwise agreed with the supervisor
    \item Two bound copies required: adhesive binding with clear cover foil and green (FAPS green) back sheet
    \item Each copy must include a CD/DVD containing:
    \begin{itemize}
        \item Digital version of the thesis (PDF)
        \item All digital sources (papers, scanned documents, websites)
        \item All images, tables, measurement data, models, and other data
    \end{itemize}
\end{itemize}

\subsection{Section and Chapter Formatting}
\begin{itemize}
    \item Chapters are numbered numerically with at most three levels (e.g., 1.1.1); no trailing dot after the last digit
    \item If a subsection X.1 exists, there must also be at least an X.2 (applies to all levels)
    \item Main chapters should have at least 3 subchapters; aim for a balanced structure (e.g., 5 chapters at level 1, each with ~4 at level 2, each with ~3 at level 3)
    \item Headings in the table of contents must match the headings in the text verbatim
    \item Every heading must be followed by actual content (including introductory text at the start of main chapters)
\end{itemize}

\subsection{Paragraphs, Spacing, and Whitespace}
\begin{itemize}
    \item Each paragraph must consist of at least two sentences
    \item No blank line between paragraphs within a chapter; one blank line before the next heading at the end of a chapter
    \item No double spaces anywhere in the document; use search-and-replace to eliminate them before submission
    \item Use non-breaking spaces (in LaTeX: \textasciitilde) between:
    \begin{itemize}
        \item Titles and names: ``Prof.~Dr.~Beispielhausen''
        \item Date expressions: ``19.~September~2007''
        \item Abbreviations: ``z.\,B.'', ``u.\,a.'', ``d.\,h.''
        \item Numbers and units: ``44~mm''
        \item Numbers and words: ``1.350~Personen'', ``Nr.~5''
        \item Citation references
    \end{itemize}
    \item Check automatic hyphenation carefully for errors (e.g., incorrect compound-word breaks)
\end{itemize}

\subsection{Figures}
\begin{itemize}
    \item Every figure receives a caption below it (Unterschrift), formulated so the figure is understandable without the surrounding text
    \item All figures are numbered consecutively with Arabic numerals and listed in the list of figures
    \item Every figure must be referenced in the body text
    \item Title image: no caption, no numbering, not included in the list of figures
    \item Appendix figures: numbered within each appendix separately (e.g., Appendix A -- Figure 1); not included in the list of figures
    \item Ensure sufficient image quality; recreate figures if originals are blurry, low-resolution, or scanned
    \item Use the same font in figures as in body text (Arial); font size 10pt is acceptable but must be consistent across all figures
    \item All figure components must be fully labeled: photo components, axes, parameters, legends with complete process parameters
    \item Fill areas with color only if the color conveys meaningful information; otherwise keep figures clean and minimal
    \item For recreated/modified figures from sources, cite with ``in Anlehnung an [xy]'' (based on [xy])
    \item Use a consistent color scheme across all figures; preferably FAPS colors unless industry cooperation requires otherwise
    \item Use a self-created graphic on the first page of the thesis
\end{itemize}

\subsection{Tables}
\begin{itemize}
    \item Every table receives a caption above it (Überschrift) with a meaningful, descriptive title
    \item Tables are numbered consecutively with Arabic numerals and listed in the table directory
    \item Every table must be referenced in the body text
    \item Use the same font as in body text; maintain consistent font size across all tables
    \item Preferred alignment inside tables: left-aligned or centered (not justified)
    \item If a table spans a page break, do not break within a row; repeat the table header on the new page
    \item Extensive tables belong in the appendix
\end{itemize}

\subsection{Diagrams}
\begin{itemize}
    \item Follow DIN 461 for axis labeling in diagrams
    \item Label axes with an arrow in the direction of increase (except for trial numbers etc.)
    \item Replace the axis label at the second-to-last tick mark with the unit
    \item Display all relevant information within the diagram
    \item Place a frame around diagrams
    \item Use Arial font inside diagrams with a legible font size
    \item Diagrams are captioned and numbered as figures (no separate numbering type)
\end{itemize}

\subsection{Formulas}
\begin{itemize}
    \item Formula numbers are enclosed in round parentheses
    \item Numbering is chapter-based in the form X.Y (X = chapter number, Y = formula number within the chapter)
    \item Numbering restarts within each main chapter
    \item Variables used in a formula must be explained directly below the formula in tabular form
\end{itemize}

\subsection{Appendix}
\begin{itemize}
    \item Technical drawings, extensive tables, supplementary documents, etc.\ go into appendices
    \item Each appendix starts on a new page and may span multiple pages
    \item Appendices are numbered with consecutive uppercase Latin letters (Appendix A, Appendix B, etc.)
\end{itemize}

\subsection{Eidesstattliche Erkl\"arung (Statutory Declaration)}
\begin{itemize}
    \item A signed declaration of independent authorship must be included on a separate page without a page number, placed after the cover page
    \item It must be signed by hand
\end{itemize}

% ============================================================================

\section{Thesis Structure and Content}

\subsection{Required Document Order}
\begin{enumerate}
    \item Deckblatt (Cover Page)
    \item Erkl\"arung des Kandidaten (Statutory Declaration)
    \item Inhaltsverzeichnis (Table of Contents)
    \item Abbildungsverzeichnis (List of Figures)
    \item Tabellenverzeichnis (List of Tables)
    \item Abk\"urzungsverzeichnis (List of Abbreviations)
    \item Einleitung und Zielstellung (Introduction and Objective)
    \item Grundlagen / Stand der Technik (Background / State of the Art)
    \item Handlungsbedarf (Research Gap / Need for Action)
    \item Methodenteil (Methodology)
    \item Ergebnisteil (Results)
    \item Diskussionsteil (Discussion)
    \item Zusammenfassung und Ausblick (Summary and Outlook)
    \item Literaturverzeichnis (Bibliography)
    \item Anhang (Appendix)
    \item Lebenslauf (Curriculum Vitae)
\end{enumerate}

\subsection{Page Count Guidelines}
\begin{itemize}
    \item Bachelorarbeit: ca.\ 60--80 pages
    \item Projekt-/Studienarbeit: ca.\ 60--80 pages
    \item Master-/Diplomarbeit: ca.\ 80--120 pages
    \item These counts exclude appendices and directories
\end{itemize}

\subsection{Red Thread (Roter Faden)}
\begin{itemize}
    \item A clear red thread must run through the entire thesis, starting from the central research question in the introduction
    \item The reader must always be able to see how any part of the thesis relates to the research question
    \item The thesis must proceed purposefully and not lose itself in tangential topics
    \item All claims and statements must be causally justified; results must not merely be listed but convincingly explained
    \item Each main chapter must begin with a short introduction motivating the content and end with a transition to the next
    \item Verify the red thread by checking all headings alone (without reading the text) --- the logical flow should be recognizable
\end{itemize}

\subsection{Chapter: Einleitung und Zielstellung (Introduction)}
\begin{itemize}
    \item Length: approximately 2 pages, maximum 3 pages
    \item Purpose: establish contact with the reader; provide orientation (``What can the reader expect?'')
    \item Content requirements:
    \begin{itemize}
        \item Introduce the topic and its general background; explain why it is relevant
        \item Present the specific research question
        \item Justify the scope and boundaries of the thesis (what is included, what is excluded, and why)
        \item Optionally describe the state of research (to show what exists and where gaps are)
        \item Explain the methodological approach
        \item Outline the structure of the thesis for the reader
    \end{itemize}
\end{itemize}

\subsection{Chapter: Grundlagen und Stand der Technik (Background and State of the Art)}
\begin{itemize}
    \item Clarify central concepts and define relevant terms
    \item Present the current state of the art objectively and without judgment (evaluation happens in the next chapter)
    \item Only cover theories, definitions, and knowledge directly relevant to the research question; avoid overly broad coverage
    \item If multiple competing theories exist, provide an overview
    \item Definitions are only needed for terms not commonly known by domain experts
    \item Present content in your own words; near-verbatim copying or translation from sources counts as insufficient
    \item Use primarily English-language sources where possible
    \item Prefer journal articles and conference papers over websites
    \item Avoid popular-science sources (e.g., Schülerduden, Brockhaus); only quote newspapers for current events
    \item Internet sources are acceptable only if no current literature exists elsewhere (e.g., software documentation); verify their credibility
\end{itemize}

\subsection{Chapter: Handlungsbedarf (Research Gap)}
\begin{itemize}
    \item Critically discuss the literature from the previous chapter regarding unsolved problems and errors
    \item Derive current deficits and research gaps; formulate the specific research question
    \item Explain who considers the problem significant and what its broader impact is
    \item Describe the contribution the thesis makes to both research (theory, models, methods, facts) and practice
    \item Narrow the problem statement into a concrete, achievable objective with clear expected outcomes
\end{itemize}

\subsection{Chapter: Methodenteil (Methodology)}
\begin{itemize}
    \item Together with Results and Discussion, this forms the main body (~2/3 of the thesis)
    \item Describe the investigation plan, procedures, and all steps and tools in sufficient detail for reproducibility
    \item Justify why specific methods were chosen over alternatives (pros/cons analysis)
    \item For experimental studies: describe the population and sample in all relevant characteristics; justify sample adequacy
    \item Scientific criteria for methods: objectivity, reliability, and validity
    \item The description must enable other researchers to reproduce the investigation
    \item Only describe analysis procedures in detail if they are novel or non-standard
\end{itemize}

\subsection{Chapter: Ergebnisteil (Results)}
\begin{itemize}
    \item Present results of the investigation objectively and without interpretation
    \item Structure the chapter logically, aligned with the theory and methodology sections
    \item Use figures, tables, and diagrams to visualize results
    \item No interpretation or evaluation at this stage
\end{itemize}

\subsection{Chapter: Diskussionsteil (Discussion)}
\begin{itemize}
    \item Summarize and critically discuss the results
    \item Evaluate and interpret results in relation to the original problem and objective from the introduction
    \item Argue the significance for scientific discourse and practical implications
\end{itemize}

\subsection{Chapter: Zusammenfassung und Ausblick (Summary and Outlook)}
\begin{itemize}
    \item Maximum 3 pages
    \item Reconnect to the central problem and objective from the introduction
    \item Compactly summarize the most important results
    \item Place results in a broader context
    \item Address:
    \begin{itemize}
        \item What are the key findings with respect to the research questions?
        \item How can the results be interpreted or evaluated?
        \item What questions remain open, and why?
        \item What do the results mean for future work on this topic?
        \item Where are opportunities for further research?
    \end{itemize}
    \item Discuss the limitations of the investigation
    \item Provide suggestions for future work and improvements based on newly emerged questions
\end{itemize}

% ============================================================================

\section{Scientific Writing Style}

\subsection{General Principles}
\begin{itemize}
    \item Write in present tense (Pr\"asens) and nominal style (Nominalstil)
    \item Use impersonal form throughout; never use ``ich'' (I) or ``man'' (one/you)
    \item Develop the red thread of argumentation before formulating sentences
    \item Every sentence and every claim must be precise and well-placed
    \item Write result-oriented: ``Die Einfl\"usse f\"uhren zu Ver\"anderungen.'' not ``Im Rahmen der Arbeit wurden untersucht\ldots''
\end{itemize}

\subsection{Sentence Construction}
\begin{itemize}
    \item Use passive voice sparingly --- prefer active formulations
    \item Avoid sentences with ``man'' entirely
    \item Be careful with infinitive constructions
    \item Avoid chaining too many nouns together (``Substantivaneinanderreihungsproblem'')
    \item No nested or overly complex sentences (Schachtels\"atze)
\end{itemize}

\subsection{Grammar}
\begin{itemize}
    \item Avoid relative time references (``k\"urzlich'', ``sp\"ater'', ``modern'') since the thesis should remain valid for years
    \item Do not use subjunctive mood (k\"onnte, sollte, m\"usste, d\"urfte)
\end{itemize}

\subsection{Technical Language}
\begin{itemize}
    \item Use precise adjectives; avoid vague terms like ``innovativ'' or ``intelligent''
    \item No pleonasms (redundant word pairs)
    \item Do not artificially inflate terms
    \item Be cautious with superlatives and comparative forms
    \item Do not invent words
    \item Use anglicisms and foreign words sparingly; only when no equivalent German term exists
    \item Use precise technical terminology; avoid colloquial language
    \item Use technical terms consistently throughout (e.g., always ``CAD-Modell'', never alternating with ``CAD Modell'')
\end{itemize}

\subsection{Expression and Phrasing}
\begin{itemize}
    \item Spell out numbers from zero to twelve; use digits from 13 onward
    \item Avoid colloquialisms and journalistic clichés
    \item Do not use rhetorical questions; make clear statements instead
    \item Avoid vague qualifiers: ``im Allgemeinen'', ``in der Regel'', ``beispielsweise'', ``teilweise''
    \item Avoid filler words: ``n\"amlich'', ``also'', ``so'', ``eben''
    \item Avoid repetition of both words and arguments
    \item Always formulate concretely and precisely; define unambiguously; enumerate completely; avoid generalizations
    \item Use nominal style: instead of ``wird bearbeitet'', write ``Zur bedarfsgerechten Bearbeitung ist auf die Einhaltung von xy zu achten.''
\end{itemize}

\subsection{Content Guidelines}
\begin{itemize}
    \item Do not document iterative development processes or suboptimal intermediate solutions; present only the final approach
    \item Do not describe device/hardware development in excessive detail; focus on scientifically relevant, generalizable, transferable aspects
    \item Use direct quotes rarely and only pointedly; prefer evaluating and developing ideas further in your own words
    \item Introduce abbreviations only when necessary for space; spell out on first use and include in the abbreviation index
    \item Emphasize scientific quality rather than engineering challenges
    \item Emphasize fundamental, generalizable, and transferable findings rather than application-specific details
    \item Instead of giving examples, enumerate completely, structure, and evaluate
    \item Make all enumerations and representations as complete and non-overlapping as possible (MECE principle)
    \item Always support claims with your own analyses or cited sources
    \item No implied personal opinion (e.g., ``jedoch'', ``\"uberaus'') and no subjective assessments except in the Outlook section
    \item Where possible, summarize each chapter's key message in a figure for quick reference
\end{itemize}

\subsection{Introductions and Transitions Between Sections}
\begin{itemize}
    \item Every main chapter or main section must begin with a brief introduction and end with an outlook/transition to the next section
    \item Example introduction: ``Hasen und Igel sind, betrachtet man ihre Laufgeschwindigkeit, auf einem vollkommen unterschiedlichen Niveau.''
    \item Example transition: ``Deshalb wird im folgenden Abschnitt der Wettlauf zwischen Hase und Igel n\"aher betrachtet.''
    \item When introducing detailed explanations, name all relevant points in the introductory sentence
    \item Bad: ``Im Folgenden werden die beiden Protagonisten n\"aher betrachtet.''
    \item Good: ``Im Folgenden werden die beiden Protagonisten Hase und Igel n\"aher betrachtet.''
\end{itemize}

\subsection{Consistency and Formatting Rules}
\begin{itemize}
    \item Introduce each abbreviation exactly once and then use it consistently
    \item Never write a paragraph consisting of only one sentence (minimum two sentences)
    \item Never write a list with only one bullet point (minimum two)
    \item Use nominal style for bullet-point items
    \item Maintain consistent spelling of expressions throughout (e.g., always ``Aerosol-Jet-Druck'', never ``Aerosoljet-Druck'')
\end{itemize}

% ============================================================================

\section{Citations and Bibliography}

\subsection{Citation Style}
\begin{itemize}
    \item Citation style: DIN ISO 690 (FAPS custom style via \texttt{bibstyleFAPS.bst})
    \item Reference sources in the text using numbers in square brackets: [xy]
    \item Page numbers are not required in citations
    \item Multiple consecutive citations use separate brackets: [xy] [yz], not [xy; yz]
    \item The bibliography is ordered chronologically and numbered consecutively; no grouping by source type
    \item Only actually cited works appear in the bibliography
\end{itemize}

\subsection{Citation Placement}
\begin{itemize}
    \item A citation before the period refers to the specific sentence [xy].
    \item A citation after the period of the last sentence refers to the entire paragraph. [xy]
\end{itemize}

\subsection{Direct Quotes}
\begin{itemize}
    \item Must be letter-perfect and character-perfect reproductions
    \item Enclose in quotation marks; cite the source at the end: ``\ldots'' [20]
    \item Quote foreign-language texts in the original language; do not translate
    \item Mark own additions or rearrangements with square brackets
    \item Mark omissions with [\ldots]
    \item Quotes longer than three lines: indent as a block, set in italic, omit quotation marks
    \item Do not correct errors in the original; mark them with [!] or [sic!]
    \item Emphasis within quotes: append [Hervorhebung durch Verfasser]
\end{itemize}

\subsection{Indirect Quotes}
\begin{itemize}
    \item Paraphrased ideas from other sources; not placed in quotation marks
    \item Still require a source reference in the same placement rules as above
\end{itemize}

\subsection{Secondary Citations}
\begin{itemize}
    \item Always prefer citing the original source
    \item Use secondary citations only when the original is unavailable; reference the secondary source in the bibliography
\end{itemize}

\subsection{Plagiarism}
\begin{itemize}
    \item Provide a source reference for every piece of knowledge that goes beyond common knowledge
    \item Originality and independence are the most important quality criteria
    \item Even when rephrasing or summarizing, the source must be cited
    \item It must be unambiguously clear to the reader at all times what is borrowed intellectual property
\end{itemize}

\subsection{Source Quality Requirements}
\begin{itemize}
    \item Every source must be: (1) citeable (publicly accessible), (2) citeworthy (scientific quality), and (3) relevant
    \item Prefer English-language sources
    \item Prefer journal articles and conference papers over websites
    \item Avoid websites as sources wherever possible; especially avoid Wikipedia, hausarbeiten.de, and similar
    \item If no conventional sources exist (e.g., for software topics), internet sources are permissible with critical evaluation
    \item Prioritize current sources from renowned journals or scientists over older work
    \item Quality indicators: scientific question, clear structure, precise terminology, theoretical foundation, appropriate methodology
\end{itemize}

\subsection{Electronic Sources}
\begin{itemize}
    \item Books, journals, and publications that also exist in print are not ``electronic sources'' even when accessed online
    \item Each source entry must answer: Who? What? Where? When?
    \item Use DOI when available; otherwise provide the full URL and access date
    \item Do not insert hyphens in URLs (could be misinterpreted as part of the address)
    \item If author is unknown: use ``o.\,V.'' (ohne Verfasser); if date is unknown: use ``o.\,J.'' or ``o.\,D.''
    \item Schema: Nachname, Vorname (Jahr): Titel. [ggf.\ weitere Angaben]. DOI oder online unter: URL (Abrufdatum). No period after DOI.
\end{itemize}

% ============================================================================

\section{Literature Research}

\begin{itemize}
    \item Recommended databases: OPAC (FAU library), Google Scholar, IEEE Digital Library, Scopus
    \item Strategy: create a search table with key terms, synonyms (abstract and specific) across multiple columns
    \item Use primarily English-language sources where possible
    \item Prefer journal articles and conference papers; avoid websites
    \item Submit a detailed outline and relevant literature to the supervisor within 2 weeks (for Abschlussarbeiten)
\end{itemize}

% ============================================================================

\section{Figures and Graphics --- Additional Rules}

\begin{itemize}
    \item Use a self-created graphic on the first page of each chapter or major section
    \item If independent creation is not possible or the graphic is assembled from multiple sources, provide a source note on the first page where it appears
    \item Poorly formatted, inconsistent, blurry, or illegible graphics create a bad impression; avoid copies and scans
    \item Recreate relevant graphics, diagrams, formulas, and tables to ensure uniform formatting and meet formal requirements
    \item Maintain a balanced ratio of text to graphics (e.g., avoid placing two graphics directly after each other)
    \item Ensure consistent design across all graphics: shapes, fonts/sizes, colors
\end{itemize}

% ============================================================================

\section{Review Checklist Before Submission}

\subsection{Content Review}
\begin{itemize}
    \item Verify all headings form a coherent red thread (recognizable without reading the text)
    \item Verify transitions and introductions exist at the beginning of every main chapter
    \item Is the central research question clearly developed in the introduction?
    \item Is the relevant literature reviewed with appropriate focus (not quantity, but quality of selection)?
    \item Can the derivation of the target concept from the analysis be followed?
    \item Are strengths and development potentials considered in the solution space?
    \item Is the proposed solution feasible and implementable?
    \item Does the outlook present innovative approaches for future work?
    \item Are all tables and figures clear and understandable?
    \item Is the context of existing literature properly established?
    \item Is there a concise summary that still covers all essential points?
    \item Is the overall structure clear and comprehensible?
    \item Is the scope sufficient yet as concise as possible?
    \item Were all key work items of the thesis topic fulfilled and the objectives achieved?
\end{itemize}

\subsection{Language and Expression Review}
\begin{itemize}
    \item Check for scientific, objective tone throughout
    \item Check for linguistic expression quality
    \item Run spell-check and grammar-check
    \item Have at least one other person proofread the thesis
\end{itemize}

\subsection{Visual and Formatting Review}
\begin{itemize}
    \item Check in single-page view for equal spacing between sections
    \item Verify balanced ratio of text and graphics (no dense clustering of figures)
    \item Verify consistent design of all graphics (shapes, fonts, sizes, colors)
    \item Eliminate all double spaces
    \item Check for sensible page breaks (no orphan headings, no broken table rows)
\end{itemize}

% ============================================================================

\section{Grading Criteria}

\begin{itemize}
    \item Systematic approach and methodology
    \item Personal effort and creativity
    \item Fulfillment of the thesis assignment
    \item Independence during the work
    \item Application of domain expertise
    \item Discussion and critical evaluation of results
    \item Presentation of the initial situation
    \item Thoroughness, completeness, and comprehensibility of the presentation
    \item External form and appearance of the thesis
    \item Expression and style
    \item Spelling and grammar
    \item Performance is measured as work divided by time --- high quality within deadlines demonstrates strong performance
\end{itemize}

% ============================================================================

\section{Presentation Guidelines (BA/PA)}

\begin{itemize}
    \item Duration: 20 minutes $\pm$ 10\% deviation
    \item Bring your own laptop; have a backup on a USB drive
    \item Submit slides to the supervisor several days before the presentation date for review
    \item Follow FAPS presentation guidelines
    \item Dress code: at minimum business casual
    \item Structure:
    \begin{enumerate}
        \item Title slide
        \item Motivation slide (between title and table of contents)
        \item Table of contents slide (all relevant topics except bibliography)
        \item Content slides
        \item Summary slide
        \item Closing slide
        \item Bibliography (hidden; only shown if asked during Q\&A). Cite sources on individual slides in (Author, Year) format
    \end{enumerate}
\end{itemize}